\chapter{Challenge-Aufgaben}

\begin{itemize}
    \item Ein Teiler $d$ einer natürlichen Zahl $n$, heisst \textit{trivial}, falls $d = 1$ oder $d = n$ gilt. Beispiele: Die trivialen Teiler von 8 sind 1 und 8 selbst. Die \textit{nichttrivialen} Teiler von 8 hingegen sind 2 und 4. Die Zahl 24 besitzt genau die sechs nichttrivialen Teiler 2, 3, 4, 6, 8 und 12 und die beiden trivialen Teiler 1 und 24.
    \item Die \textit{Quersumme} einer Zahl ist die Summe der einzelnen Ziffern ihrer Darstellung im Dezimalsystem. Beispiele: Die Quersumme von 315 ist $3 + 1 + 5 = 9$. Die Quersumme von 1004 ist $1 + 4 = 5$.
    \item Anstelle der Quersumme können wir auch die \textit{Quersumme der quadrierten Ziffern} berechnen, indem wir die Ziffern jeweils zuerst quadrieren, bevor wir sie summieren. Beispiel: Die Quersumme der quadrierten Ziffern von 52 ist $5^2 + 2^2 = 25 + 4 = 29$. Die Quersumme der quadrierten Ziffern von 134 ist $1^2 + 3^2 + 4^2 = 1 + 9 + 16 = 26$.
    \item Eine natürliche Zahl $n$, welche keine Primzahl ist, heisst \tib{kanada-Perfekt}, falls die Summe der nichttrivialen Teiler von $n$ der Quersumme der quadrierten Ziffern von $n$ entspricht.
    \item Die kleinste kanada-perfekte Zahl ist 125. Die nichttrivialen Teiler von 125 sind 5 und 25. Somit ist die Summe der nichtrivialen Teiler von 125 gegeben durch $5 + 25 = 30$. Die Quersumme der quadrierten Ziffern von 125 ist ebenfalls 30 , denn $1^2 + 2^2 + 5^2 = 30$. Somit ist 125 eine kanada-perfekte Zahl.
\end{itemize}

\noindent
Schreiben Sie ein Python-Programm, welches alle kanada-perfekten Zahlen kleiner als 600 findet. Hinweis: Es gibt genau zwei.
